\chapter{Summary and Outlook}

This thesis presents the realization of an atom-photon quantum gate between a single \textsuperscript{87}Rb atom trapped in a high-finesse birefringent Fabry-Pérot cavity and a polarization-encoded photonic qubit, realized experimentally as a weak coherent pulse. This represents a further step towards a fully functional quantum network node on the CavityX platform, complementing the previously shown heralded atom-photon entanglement \cite{HeraldedEntanglement} and quantum memory capabilities \cite{QuantumNetworkManuel} of the setup. \\

The gate mechanism exploits the strong cavity birefringence, which splits the polarization eigenmodes by $\Delta_\text{biref} \approx 2\pi \times \SI{505}{\mega\hertz}$, and the atom-state-dependent reflection of the $\pi$-polarized mode. An incoming $\pi$ polarized photon acquires a $\pi$ phase shift when the atom is prepared in the uncoupled clock state $\ket{0} = \ket{F{=}1, m_F{=}0}$, but no phase shift when the atom is prepared in the coupled clock state $\ket{1} = \ket{F{=}2, m_F{=}0}$. Using this I have implemented a controlled-phase-flip gate with the photon being encoded in the $\{\ket{H}, \ket{V}\}$ basis, where $\ket{H}$ is the polarization that drives the $\pi$ atomic transition and will henceforth be denoted $\ket{\pi}$. By further preparing the atom in the superposition basis $\{\pm\equiv\frac{\ket{V}\pm\ket{\pi}}{\sqrt{2}}\}$, I was able to realize a controlled NOT gate using the same working principles. Because the cavity's $V$-mode is off-resonant, a photon with such polarization is reflected with near-unity amplitude. On the other hand, the $\pi$-mode suffers intracavity losses as its reflection amplitude is determined by the atomic qubit state. This imbalance in reflection amplitudes between the two polarizations would lead to a loss in overlap with the ideal gate's truth table. Thus, a custom mount holding an array of UV fused silica Brewster windows was installed in the input path to attenuate the $V$-polarized component externally and match the two reflection amplitudes.\\

The cavity and atom-cavity parameters were characterized experimentally, yielding an atom-cavity coupling strength of $g = 2\pi \times (25.9 \pm 0.5)\,\text{MHz}$. A master-equation simulation was used to determine the optimal mean input photon number $\hat{n} \approx 0.1$ and $V$-mode transmission $T_V \approx \SI{24}{\percent}$ (corresponding to six to seven Brewster windows) for gate operation.\\

Using weak coherent pulses with a mean input photon number of $\bar{n}\approx0.1$ and a configuration with seven Brewster windows installed in the custom Brewster mount, overlaps of $(92.41 \pm 0.52)\,\%$ with the ideal CPF truth table and $(88.15 \pm 0.86)\,\%$ with the ideal CNOT truth table were achieved. The coherent action of the gate was further demonstrated by generating an atom-photon Bell state $\ket{\Phi^+}$ from a separable input, reaching a fidelity of $(65.90 \pm 1.73)\,\%$ -- above the \SI{50}{\percent} threshold required to confirm entanglement generation. The lower Bell-state fidelity compared to the truth-table overlap reflects the additional sensitivity of $F_{\text{Bell}}$ to phase errors, which do not contribute to the classical truth-table overlap. A detailed attribution of the infidelity to specific error sources is left to future work. Scans of the overlap against both the $V$-mode attenuation and the mean photon number agreed with the master-equation simulation, apart from an enhanced drop at $\bar{n} \gtrsim 1$ whose origin remains to be identified. The dominant contributions to the gate infidelity were identified as the asymmetry between the coupled and uncoupled $\pi$-mode reflectivities ($\sim5\%$) and the imperfect atomic state preparation of the coupled atomic state ($\sim6\%$). The former is set by the current cavity cooperativity and imposes a ceiling on the achievable overlap for the given cavity parameters, whereas the latter can be reduced by refining the individual components of the state-preparation sequence. Smaller contributions that limit the gate overlap include thermal drifts of the detection basis and run-to-run fluctuations of the atom-cavity coupling.\\

With the truth tables confirming the correct input-output behavior and the generation of a Bell state confirming the coherent action of the gate, the natural next step is to improve the Bell state fidelity through a more systematic optimization of the experimental parameters -- in particular, the atomic state preparation fidelity and the photon polarization contrast. Beyond that, inter-node protocols such as quantum state teleportation via a fiber link between the MPQ in Garching and the LMU in Munich come within reach~\cite{TelecomConnectionMPQLMU}. Another compelling direction is the investigation of nonlinear photon-photon interactions mediated by the single atom, as in the photon-photon gate proposed by Duan and Kimble \cite{DuanKimble2004}. A related scheme enabled by the existing infrastructure is the generation of photon-photon entanglement by combining the existing quantum memory capability with the gate demonstrated here: a first photon incident via the long cavity is mapped onto the atomic state, and a second photon is subsequently reflected off the atom-cavity system, producing entanglement between the two temporally separated photons. In contrast to the sequential two-photon reflection scheme of Hacker \textit{et al.}~\cite{Hacker2016}, the storage of the first photon provides a heralding signal that confirms its successful arrival before the second photon is sent, allowing the protocol to be conditioned on a verified first photon.
